% !TEX TS-program = pdflatex
% !TeX program = pdflatex
% !TEX encoding = UTF-8
% !TEX spellcheck = fr

\documentclass[xcolor=table]{beamer}


%\usepackage{fullpage}
%\usepackage[left=2.8cm,right=2.2cm,top=2 cm,bottom=2 cm]{geometry}
\setbeamersize{text margin left=10pt,text margin right=10pt}
\usepackage{amsmath,amssymb} 
\usepackage[T1]{fontenc}
\usepackage{textcomp}
\usepackage[utf8]{inputenc}
\usepackage[french, onelanguage]{algorithm2e}
\usepackage[french]{babel}
\usepackage{arabtex}
\usepackage{txfonts}
\usepackage[]{graphicx}
\usepackage{multirow}
\usepackage{hyperref}
\usepackage{colortbl}
\usepackage{listingsutf8}
\usepackage{wrapfig}
\usepackage{multicol}
\usepackage[export]{adjustbox} %for images in table, also for frame

\hypersetup{
	colorlinks,
	urlcolor = blue
}

%\renewcommand{\baselinestretch}{1.5}

\def\supit#1{\raisebox{0.8ex}{\small\it #1}\hspace{0.05em}}

\AtBeginSection{%
	\begin{frame}
		\sectionpage
	\end{frame}
}

\newcommand{\rottext}[2]{%
	\rotatebox{90}{%
	\begin{minipage}{#1}%
		\raggedleft#2%
	\end{minipage}%
	}%
}

\usepackage{longtable}
\usepackage{tabu}


\institute{ %
École  nationale Supérieure d'Informatique (ESI, ex. INI), Algérie
}
\author[ \textbf{\footnotesize\insertframenumber/\inserttotalframenumber} \hspace*{1.5cm} ESI - BENATCHBA \& ARIES (2019-2020)] %
{BENACHBA Karima \hfill ARIES Abdelkrime}

\titlegraphic{\includegraphics[height=1cm]{assets/esi-logo.png}}%\hspace*{4.75cm}~

\date{Année universitaire : 2019/2020} %\today

\usetheme{Warsaw} % Antibes Boadilla Warsaw

\beamertemplatenavigationsymbolsempty

%\setbeamertemplate{headline}{}

\definecolor{lightblue}{HTML}{D0D2FF}
\definecolor{lightyellow}{HTML}{FFFFAA}
\definecolor{darkblue}{HTML}{0000BB}
\definecolor{olivegreen}{HTML}{006600}
\definecolor{violet}{HTML}{6600CC}

\newcommand{\keyword}[1]{\textcolor{red}{\bfseries\itshape #1}}
\newcommand{\expword}[1]{\textcolor{olivegreen}{#1}}
\newcommand{\optword}[1]{\textcolor{violet}{\bfseries #1}}

\makeatletter
\newcommand\mysphere{%
	\parbox[t]{10pt}{\raisebox{0.2pt}{\beamer@usesphere{item projected}{bigsphere}}}}
\makeatother

%\let\oldtabular\tabular
%\let\endoldtabular\endtabular
%\renewenvironment{tabular}{\rowcolors{2}{white}{lightblue}\oldtabular\rowcolor{blue}}{\endoldtabular}


\NoAutoSpacing %french autospacing after ":"

\def\graphpath{}

\newcommand{\changegraphpath}[1]{\def\graphpath{#1}}


\newcommand{\vgraphpage}[2][.84\textheight]{%
%	\begin{center}%
		\includegraphics[height=#1]{\graphpath #2}%
%	\end{center}%
}

\newcommand{\hgraphpage}[2][\textwidth]{%
%	\begin{center}%
		\includegraphics[width=#1]{\graphpath #2}%
%	\end{center}%
}

\newcommand{\graphpage}[2][]{%
	\includegraphics[#1]{\graphpath #2}%
}

\bibliographystyle{acm}

\newcommand{\insertbibliography}[2]{
	\appendix
	\section*{Bibliographie}
	\nocite{#2}
%	\makeatletter % to change template
%	\setbeamertemplate{headline}[default] % not mandatory, but I though it was better to set it blank
%	\def\beamer@entrycode{\vspace*{-\headheight}} % here is the part we are interested in :)
%	\makeatother
	\begin{multicols*}{2}[\frametitle{\insertsection} \usebeamertemplate*{frametitle}]%\usebeamertemplate*{frametitle}\frametitle{Références}
		\tiny
		\bibliography{#1}
	\end{multicols*}
}

\newcommand{\insertlicence}{
	\begin{frame}[plain]
	\frametitle{Licence : CC-BY-NC 4.0}
%	\framesubtitle{Licence: CC-BY-NC 4.0}
	\begin{center}
		\bfseries Attribution - Pas d'Utilisation Commerciale 4.0 International 
		\url{https://creativecommons.org/licenses/by-nc/4.0/deed.fr}
	\end{center}

	\tiny

	Vous êtes autorisé à : 
	\begin{itemize}
		\item \textbf{Partager} — copier, distribuer et communiquer le matériel par tous moyens et sous tous formats
		\item \textbf{Adapter} — remixer, transformer et créer à partir du matériel
	\end{itemize}
	
	Selon les conditions suivantes : 
	\begin{itemize}
		\item \textbf{Attribution} — Vous devez créditer l'Œuvre, intégrer un lien vers la licence et indiquer si des modifications ont été effectuées à l'Oeuvre. Vous devez indiquer ces informations par tous les moyens raisonnables, sans toutefois suggérer que l'Offrant vous soutient ou soutient la façon dont vous avez utilisé son Oeuvre.
		\item \textbf{Pas d'Utilisation Commerciale} — Vous n'êtes pas autorisé à faire un usage commercial de cette Oeuvre, tout ou partie du matériel la composant. 
		\item \textbf{Pas de restrictions complémentaires} — Vous n'êtes pas autorisé à appliquer des conditions légales ou des mesures techniques qui restreindraient légalement autrui à utiliser l'Oeuvre dans les conditions décrites par la licence.
	\end{itemize}

	\end{frame}
}

\settowidth{\leftmargini}{\usebeamertemplate{itemize item}}
\addtolength{\leftmargini}{\labelsep}

\AtBeginDocument{
	\newcolumntype{L}[2]{>{\vbox to #2\bgroup\vfill\flushleft}p{#1}<{\egroup}} 
	
	\begin{frame}[plain]
		\maketitle
	\end{frame}

%	\insertlicence
}


% needs etoolbox; to break links after -
\appto\UrlBreaks{\do\-}


\title[ML: Naïve Bayes] %
{Machine Learning \\Naïve Bayes}  

\changegraphpath{../img/naive-bayes/}

\begin{document}

\begin{frame}
\frametitle{Naïve Bayes}
\framesubtitle{Rappel}
\huge
\[
\overbrace{P(A|B)}^\text{Posérieure} = \frac{\overbrace{P(A)}^\text{Antérieure} \overbrace{P(B|A)}^{\text{Vraissemblance}}}{\underbrace{P(B)}_\text{Evidence}}
\]
\end{frame}

\begin{frame}
\frametitle{Naïve Bayes}
\framesubtitle{Plan}

\begin{multicols}{2}
	%	\small
	\tableofcontents
\end{multicols}
\end{frame}

\section{Modèles de classification}

\subsection{Modèle discriminatif}

\begin{frame}
	\frametitle{Naïve Bayes}
	\framesubtitle{Modèles de classification : modèle discriminatif}
	
	\begin{itemize}
		\item Le modèle apprend la probabilité conditionnelle $P(Y|X=x)$
		\item apprendre les frontières entre les classes
		\item Exemples des classificateurs 
		\begin{itemize}
			\item Régression logistique
			\item SVM
			\item Arbres de décision 
			\item Réseau de neurones 
			\item ...
		\end{itemize}
	\end{itemize}
	
\end{frame}

\subsection{Modèle génératif}

\begin{frame}
	\frametitle{Naïve Bayes}
	\framesubtitle{Modèles de classification : modèle génératif}
	
	\begin{itemize}
		\item Le modèle apprend la probabilité conditionnelle $P(X|Y=y)$
		\item apprendre la distribution des classes individuelles
		\item Exemples des classificateurs 
		\begin{itemize}
			\item Naïve Bayes
			\item Modèle de Markov caché
			\item Auto-encodeur variationnel
			\item ...
		\end{itemize}
	\end{itemize}
	
\end{frame}

\section{Théorie}

\subsection{Classification}

\begin{frame}
	\frametitle{Naïve Bayes}
	\framesubtitle{Théorie : Classification}
	
	\[P(Y_k|X) = \frac{P(Y_k) P(X|Y_k)}{X}\]
	
	\[P(Y_k|X) \propto P(Y_k) P(X|Y_k)\]
	
	\[P(Y_k|X) \propto P(Y_k) \prod_{j=1}^{N} P(X_j|Y_k)\]
	
	\[\hat{Y} = \arg\max_{Y_k} P(Y_k) \prod_{j=1}^{N} P(X_j|Y_k)\]
	
	\[\hat{Y} = \arg\max_{Y_k} \log P(Y_k) + \sum_{j=1}^{N} \log P(X_j|Y_k)\]
	
\end{frame}

\subsection{Apprentissage}

\begin{frame}
	\frametitle{Naïve Bayes}
	\framesubtitle{Théorie : Apprentissage (probabilité antérieure)}
	
	\[P(Y_k) = \frac{|\{y / y \in Y \wedge y = k\}|}{M}\]
	
	\begin{itemize}
		\item $|\{y / y \in Y \wedge y = k\}|$ est le nombre des échantillons appartenant à la classe $k$ dans le dataset d'entraînement
		\item $M$ est le nombre des échantillons dans le dataset d'entraînement
		\item Si distribution des échantillons d'entraînement est uniforme sur les classes, cette probabilité est négligées 
		\item Si on veut donner le même poids aux classes, on néglige cette probabilité
	\end{itemize}
	
	
\end{frame}


\begin{frame}
	\frametitle{Naïve Bayes}
	\framesubtitle{Théorie : Apprentissage (vraisemblance - loi multinomiale)}
	
	\[P(x_j = v_j|Y_k) = \frac{|X_j = v_j|_{Y_k}}{\sum_{k'=1}^{L} |X_j = v_j|_{Y_k'}}\]
	
	\begin{itemize}
		\item $x_j$ est la caractéristique $j$
		\item $v_j$ est une valeur parmi les valeurs existantes dans la caractéristique $j$
		\item $|X_j = v_j|_{Y_k} = |\{Y^{(i)} / Y^{(i)} = k \wedge X^{(i)}_j = v_j\}|$ le nombre des échantillons de la caractéristique $j$ ayant la valeur $v_j$ et appartenant à la classe $k$
		\item $V_j$ est la liste des valeurs uniques dans la caractéristique $j$
		\item \expword{sklearn.naive\_bayes.CategoricalNB}
	\end{itemize}

	\[P(x_j = v_j|Y_k) = \frac{|X_j = v_j|_{Y_k} + \alpha}{\sum_{k'=1}^{L} |X_j = v_j|_{Y_k'} + \alpha |V_j|}\]
	
	
\end{frame}

\begin{frame}
	\frametitle{Naïve Bayes}
	\framesubtitle{Théorie : Apprentissage (vraisemblance - loi multinomiale - texte)}
	
	\begin{itemize}
		\item Un texte peut être vu comme une seule caractéristique où les mots sont les catégories
		\item $C_j$ est le compte du mot $w_j$ dans tous les textes
		\item $C_{jk}$ est le compte du mot $w_j$ dans les textes appartenant à la classe $Y_k$
		\item $V$ est le vocabulaire (les différents mots dans le dataset d'entraînement)
		\item \expword{sklearn.naive\_bayes.MultinomialNB}
	\end{itemize}
	
	\[P(mot = w_j|Y_k) = \frac{C_{jk} + \alpha}{C_j + \alpha |V|}\]
	\[\hat{Y} = \arg\max_{Y_k} P(Y_k) * \prod_{m \in texte} P(mot = m|Y_k)\]
	
\end{frame}

\begin{frame}
	\frametitle{Naïve Bayes}
	\framesubtitle{Théorie : Apprentissage (vraisemblance - loi de Bernoulli)}
	
	\[P(x_j|Y_k) = P(X_j|Y_k) x_j + (1-P(X_j|Y_k)) (1-x_j)\]
	\[P(X_j|Y_k) = \frac{|\{X_j^{(i)} = 1 \wedge y^{(i)} = k\}|}{|\{y^{(i)} = k\}|}\]
	\begin{itemize}
		\item $x_j \in \{0, 1\}$ est une valeur donnée de la caractéristique $j$
		\item $X_j$ est l'ensemble des valeurs booléennes de l'attribut $j$ dans l'entraînement
		\item \expword{sklearn.naive\_bayes.BernoulliNB}
	\end{itemize}
	
	
\end{frame}

\begin{frame}
	\frametitle{Naïve Bayes}
	\framesubtitle{Théorie : Apprentissage (vraisemblance - loi normale)}
	
	\[P(X_j = v|Y_k) = \frac{1}{2\pi \sigma_{kj}^2} e^\frac{-(v-\mu_{kj})^2}{2 \sigma_{kj}^2}\]
	\begin{itemize}
		\item $x_j \in \{0, 1\}$ est représentée par une valeur numérique
		\item $\mu_{kj}$ est l'espérance (la moyenne) des valeurs de l'attribut $j$ appartenant à la classe $k$
		\item $\sigma_{kj}$ est la variance des valeurs de l'attribut $j$ appartenant à la classe $k$
		\item $v$ est une valeur donnée
		\item \expword{sklearn.naive\_bayes.GaussianNB}
	\end{itemize}
	
	
\end{frame}

\section{Application}

\subsection{Loi multinomiale}

\begin{frame}
	\frametitle{Naïve Bayes}
	\framesubtitle{Application : Loi multinomiale}
	
	\scriptsize
	\rowcolors{2}{lightblue}{lightyellow}
	\begin{center}
		\begin{tabular}{lllll}
		\rowcolor{darkblue}
		\textcolor{white}{temps} & \textcolor{white}{temperature} & \textcolor{white}{humidite} & \textcolor{white}{vent} & \textcolor{white}{jouer} \\
		ensoleile & chaude & haute & non & non \\
		ensoleile & chaude & haute & oui & non \\
		nuageux & chaude & haute & non & oui \\
		pluvieux & douce & haute & non & oui \\
		pluvieux & fraiche & normale & non & oui \\
		pluvieux & fraiche & normale & oui & non \\
		nuageux & fraiche & normale & oui & oui \\
		ensoleile & douce & haute & non & non \\
		ensoleile & fraiche & normale & non & oui \\
		pluvieux & douce & normale & non & oui \\
		ensoleile & douce & normale & oui & oui \\
		nuageux & douce & haute & oui & oui \\
		nuageux & chaude & normale & non & oui \\
		pluvieux & douce & haute & oui & non \\
	\end{tabular}
	\end{center}

	\begin{itemize}
		\item Trouver la classe de \{pluvieux, chaude, haute, non \}
	\end{itemize}
	
	
\end{frame}


\subsection{Loi de Bernoulli}

\begin{frame}
	\frametitle{Naïve Bayes}
	\framesubtitle{Application : Loi de Bernoulli}
	
	\scriptsize
	\rowcolors{2}{lightblue}{lightyellow}
	\begin{center}
		\begin{tabular}{llll}
			\rowcolor{darkblue}
			\textcolor{white}{confiant} & \textcolor{white}{étudié} & \textcolor{white}{malade} & \textcolor{white}{résultat} \\
			Oui & Non & Non & échouer \\
			Oui & Non & Oui & passer \\
			Non & Oui & Oui & échouer \\
			Non & Oui & Non & passer \\
			Oui & Oui & Oui & passer \\
		\end{tabular}
	\end{center}
	
	\begin{itemize}
		\item Trouver la classe de \{Non, Non, Non \}
	\end{itemize}
	
	
\end{frame}


\subsection{Loi normale}

\begin{frame}
	\frametitle{Naïve Bayes}
	\framesubtitle{Application : Loi normale}
	
	\scriptsize
	\rowcolors{2}{lightblue}{lightyellow}
	\begin{center}
		\begin{tabular}{llll}
			\rowcolor{darkblue}
			\textcolor{white}{taille} & \textcolor{white}{poids} & \textcolor{white}{pointure} & \textcolor{white}{sexe} \\
			182 & 81.6 & 30 & masculin\\
			180 & 86.2 & 28 & masculin\\
			170 & 77.1 & 30 & masculin\\
			180 & 74.8 & 25 & masculin\\
			152 & 45.4 & 15 & féminin\\
			168 & 68.0 & 20 & féminin\\
			165 & 59.0 & 18 & féminin\\
			175 & 68.0 & 23 & féminin\\
		\end{tabular}
	\end{center}
	
	\begin{itemize}
		\item Trouver la classe de \{183, 59, 20 \}
	\end{itemize}
	
	
\end{frame}

%\insertbibliography{ML-RN}{*}

\end{document}

